\section{Esperienze accademiche}
%
\outerlist{

\entrybig
	{\textbf{Ricercatore a Tempo Determinato di Tipo A (RTDa)}}{AimageLab, UNIMORE}
	{}{In corso}
 \innerlist{
	\entry{\textbf{ECOSISTER (progetto PNRR, \href{https://ecosister.it/spoke/spoke-4/}{Spoke 4})}: “Pedestrian and cyclist safety, high-quality cycling network, modelling mobility flows, multimodal systems and shared mobility, cybernetic mobility, intelligent video system”.}
    \entry{\textbf{Attività:} nell’ambito dello \textbf{Spoke 4} del progetto ECOSISTER, le attività di ricerca hanno riguardato lo studio, la progettazione e lo sviluppo di sistemi di \textit{Artificial Intelligence} e \textit{Computer Vision} a supporto di applicazioni di \textit{smart city}, con particolare enfasi sulla percezione e l’analisi della densità del traffico pedonale in contesti urbani. Le attività svolte si collegano alle linee di ricerca in \textit{Multiple Object Tracking}, \textit{People Detection and Tracking}, \textit{Trajectory Prediction} e \textit{Continual Learning}.}
}

\entrybig
	{\textbf{Assegnista di ricerca}}{AimageLab, UNIMORE}
	{Advisor: Prof. Simone Calderara}{2022}
\innerlist{
	\entry{\textbf{Attività:} nell’ambito del progetto europeo Horizon 2020 \href{https://www.insectt.eu}{\textbf{InSecTT}} (\textit{Intelligent Secure Trustable Things}), studio, progettazione e sviluppo di tecniche di \textit{Deep Learning} per l’individuazione di comportamenti e/o eventi anomali nel trasporto pubblico. Contestualmente, sono state svolte attività di coordinamento locale delle attività di sviluppo, nonché la redazione di documentazione tecnica e scientifica (es. report, presentazioni) a supporto delle fasi di revisione e rendicontazione del progetto.}
}

\entrybig
	{\textbf{Dottorato di ricerca}}{AimageLab, UNIMORE}
	{Advisor: Prof. Simone Calderara}{2019 -- 2021}
\innerlist{
    \entry{\textbf{Temi di ricerca trattati:} anomaly detection, apprendimento continuo, riconoscimento facciale (in ambito umano, animale e veicolare), analisi di immagini satellitari mediante reti neurali, self-supervised learning.}
    \entry{\textbf{Attività:} studio e analisi dello stato dell’arte, progettazione e sviluppo di modelli di \textit{Deep Learning}. Scrittura di articoli scientifici, gestione del processo di revisione e successiva pubblicazione. Presentazione dei risultati della ricerca presso conferenze internazionali. Attività di supporto alla didattica e di advisory per tesi di laurea magistrale.}
}

}